\subsection{Integrimi Monte Carlo si vlerësim i një pritjeje}
Le të jetë $X\sim p$ dhe
\begin{equation}
 I=\int h(x)p(x)\dd x=\E_p[h(X)].
\end{equation}
Vlerësuesi bazë është
\begin{equation}
 \widehat I_N=\frac1N\sum_{i=1}^{N}h(X_i).
\end{equation}
Për mostra të pavarura,
\begin{equation}
 \E[\widehat I_N]=I,
 \qquad
 \Var(\widehat I_N)=\frac{\sigma_h^2}{N}.
\end{equation}
Gabimi standard vlerësohet me $s_h/\sqrt N$. Teorema qendrore kufitare jep
\begin{equation}
 \frac{\widehat I_N-I}{s_h/\sqrt N}\approx\mathcal N(0,1).
\end{equation}
Kështu për ta përgjysmuar gabimin duhen afërsisht katër herë më shumë mostra. Shkalla $N^{-1/2}$ nuk varet drejtpërdrejt nga dimensioni, por varianca e integrandit dhe vështirësia e kampionimit mund të rriten fort me dimensionin.

Për integralin $J=\int_D f(\vect x)\dd\vect x$ mbi një rajon me vëllim $V_D$, nëse $X$ është uniforme në $D$, atëherë $J=V_D\E[f(X)]$. Metoda gjeometrike hit-or-miss përdor një ndryshore treguese dhe zakonisht ka variancë më të madhe se mesatarja e drejtpërdrejtë e $f$.

\subsection{Variablat antitetike}
Nëse $U\sim\mathcal U(0,1)$, edhe $1-U$ është uniforme. Vlerësuesi
\begin{equation}
 \widehat I_{\mathrm{anti}}=\frac1N\sum_{i=1}^{N}
 \frac{h(U_i)+h(1-U_i)}{2}
\end{equation}
ka variancë
\begin{equation}
 \frac{1}{2N}\left[\Var(h(U))+\Cov(h(U),h(1-U))\right].
\end{equation}
Për $h$ monotone, kovarianca zakonisht është negative dhe varianca ulet. Metoda nuk është automatikisht e dobishme për funksione jo monotone; kovarianca duhet matur.

\subsection{Variabli kontrollues}
Le të jetë $Y$ madhësia e interesit dhe $C$ një madhësi e korreluar me pritje të njohur $\mu_C$. Vlerësuesi
\begin{equation}
 Y_\beta=Y-\beta(C-\mu_C)
\end{equation}
mbetet i paanshëm. Varianca është
\begin{equation}
 \Var(Y_\beta)=\Var(Y)+\beta^2\Var(C)-2\beta\Cov(Y,C).
\end{equation}
Derivimi ndaj $\beta$ jep optimumin
\begin{equation}
 \beta_*=\frac{\Cov(Y,C)}{\Var(C)}.
\end{equation}
Në praktikë $\beta_*$ vlerësohet nga një mostër pilot ose nga e njëjta mostër, duke raportuar qartë procedurën.

\subsection{Kampionimi i rëndësisë për ngjarje të rralla}
Kur integrali dominohet nga një rajon që $p$ e viziton rrallë, kampionimi i drejtpërdrejtë prodhon shumicën e kontributeve zero dhe disa pesha shumë të mëdha. Me $X\sim q$,
\begin{equation}
 I=\E_q\left[h(X)\frac{p(X)}{q(X)}\right].
\end{equation}
Propozimi ideal formal është $q_*(x)\propto|h(x)|p(x)$, por varet nga integrali i panjohur. Ai shërben si parim projektimi: $q$ duhet të imitojë madhësinë e integrandit, jo vetëm dendësinë bazë.

\subsection{Përhapja Monte Carlo në zbritjen e mjetit hapësinor}
Në një model të thjeshtuar vertikal,
\begin{align}
 \dot h&=v,\\
 m\dot v&=-mg(h)+\tfrac12\rho(h)C_DA,v|v|+T(t),
\end{align}
ku $h$ është lartësia dhe $v<0$ gjatë zbritjes. Parametrat $C_D$, masa, dendësia atmosferike dhe shtytja kanë pacaktueshmëri. Çdo realizim përdor një mostër të parametrave, integron ODE-në dhe regjistron shpejtësinë në kontakt. Probabiliteti i dështimit vlerësohet nga treguesi $I_r=1$ kur $|v_{\mathrm{kontakt}}|$ tejkalon kufirin.

Një analizë bindëse raporton intervalin e besimit të probabilitetit, jo vetëm numrin e dështimeve. Për ngjarje shumë të rralla, simulimi i drejtpërdrejtë është joefikas; nevojitet kampionim i rëndësisë ose ndarje e ngjarjeve, me korrigjimin përkatës të peshave.

\begin{lstlisting}[language=Python,caption={Vlerësues Monte Carlo me gabim standard.}]
import numpy as np

def monte_carlo(h, gjenero, N, fare=2026):
    rng = np.random.default_rng(fare)
    x = gjenero(rng, N)
    y = np.asarray(h(x), float)
    vleresimi = y.mean()
    gabimi_standard = y.std(ddof=1) / np.sqrt(N)
    return vleresimi, gabimi_standard
\end{lstlisting}

\subsection{Përmbledhje dhe ushtrime}
Monte Carlo prodhon një vlerësim dhe një procedurë të matjes së gabimit. Numri i shifrave duhet të jetë në përputhje me gabimin standard.
\begin{enumerate}
\item Nxirrni variancën e metodës hit-or-miss për një sipërfaqe me probabilitet $p$.
\item Provoni formulën e $\beta_*$ dhe zvogëlimin maksimal të variancës në terma të korrelacionit.
\item Krahasoni kuadraturën deterministe dhe Monte Carlo kur dimensioni rritet.
\item Ndërtoni një analizë konvergjence me shumë fara dhe përshtatni pjerrësinë në grafik log--log.
\item Vlerësoni probabilitetin e uljes së pasigurt të mjetit dhe jepni interval Wilson.
\end{enumerate}
